\chapter{FPGA տվյալների գեներացիա և մշակում}
\label{ch:technical}

Սկզբնական տարբերակում իրական սենսորները պարտադիր չեն։ Միկրոէլեկտրոնիկայի
ուսանողը պետք է իրականացնի \textbf{կոնֆիգուրելի ազդանշան-գեներատոր / սենսոր-էմուլյատոր}
RTL մոդուլներ, որոնք FPGA-ի ներսում մոդելավորում են սենսորի իրատեսական վարքը։
Հնարավոր սենսորների տիպերը՝ ջերմաստիճան (Temperature), ճնշում (Pressure),
վիբրացիա (Vibration), ընդհանրական անալոգային սենսոր և այլ կոնֆիգուրելի մոդելներ։

Յուրաքանչյուր սենսորի համար առանձին ապարատ իրականացնելու փոխարեն նախատեսվում են
\textbf{վերաօգտագործելի ազդանշան-գեներացման պրիմիտիվներ}.
\begin{itemize}
  \item \textbf{Ramp Generator} — աստիճանաբար աճող/նվազող արժեքներ (ջերմաստիճան, ճնշում)։
  \item \textbf{Periodic Signal Generator} — պարբերական ազդանշաններ (վիբրացիա)։
  \item \textbf{Pseudo-Random / Noise Generator} — \gls{lfsr}-ի հիման վրա, կոնֆիգուրելի աղմուկ։
  \item \textbf{Anomaly / Fault Generator} — ոչ նորմալ վարքի ներարկում (գերտաքացում, բարձր վիբրացիա)։
\end{itemize}

\adddiagram[fig:rtl]{rtl-pipeline}{0.32}{FPGA-ի RTL մշակման հնարավոր փոխակրիչը (pipeline).}

\textbf{Signal Generator}-ի կոնֆիգուրելի պարամետրերն են՝ ազդանշանի տիպ, ամպլիտուդ,
հաճախություն, բազային արժեք, նմուշառման արագություն, աղմուկի մակարդակ, անոմալիայի
միացում և մեծություն։ \textbf{Sensor Emulator}-ը ընդհանրական ազդանշանը վերածում է
սենսորանման արժեքների՝ scaling, offset, saturation, աղմուկի ու անոմալիայի ներարկման միջոցով։

FPGA-ն կատարում է իրական ժամանակի մշակում (\textbf{Processing IP}-ներ). հնարավոր
մոդուլներ՝ Moving Average, Min/Max Detector, Threshold Detector, Event Detector,
\gls{fir} Filter, Statistical accumulator։ MVP-ն պարունակում է \textbf{3 մշակող մոդուլ}.
որոնք ընտրել՝ որոշում է ուսանող 1-ը ղեկավարի հետ։
